\chapter*{Abstract}
\addcontentsline{toc}{chapter}{Abstract}

Increasingly heterogeneous hardware, including CPUs, GPUs, and accelerators, often relies on shared-memory interfaces to access and exchange data. However, different architectures can allow different shared-memory behaviors, so the same concurrent program may behave differently across hardware. Memory consistency models describe these behaviors by specifying which memory values each thread is allowed to observe. To reason about these behaviors reliably, we need precise descriptions of memory consistency models. Empirical testing can expose some behaviors, but it cannot characterize all possible executions. Formal reasoning provides a more precise foundation, but pen-and-paper proofs are difficult to check, reuse, and maintain.

Mechanization addresses these limitations by encoding memory consistency model definitions and proofs in a proof assistant, where they can be checked by a computer. However, building mechanized developments requires additional effort. We present \textsc{LeanCat}, a domain-specific proof framework for mechanized reasoning about memory consistency models in Lean~4. \textsc{LeanCat} builds on Mathlib and existing specifications of memory consistency models, providing reusable definitions, abstractions, and proved lemmas for reasoning about weak memory consistency models. It supports applications such as comparing the strength of different memory consistency models, proving instruction-mapping correctness between architectures, and designing memory consistency models for systems that combine devices with different memory consistency models. \textsc{LeanCat} therefore provides a practical foundation for verified reasoning about weak memory consistency models, with implications for compiler validation and cross-platform correctness guarantees.
